%! Operations with Polynomials — Unit 4, Lesson 4.1 — Group Activity
\documentclass[10pt]{article}
\usepackage{algebra2-article}
\usepackage{algebra2-boxes}
\newcommand{\TallMath}[1]{$\displaystyle #1\rule[-1.4em]{0pt}{3.2em}$}

\begin{document}

\pageheader{Unit 4, Lesson 4.1}{Group Activity}
\namepartnerperiod

\vspace{0.08in}

\begin{headlinebox}{forestbg}
\small\textbf{Building Polynomials.} Yesterday you \emph{read} polynomials off their graphs; today you
\emph{build} them. Two rules carry everything: \emph{like terms combine, unlike terms do not}, and
\emph{every term times every term}. Write every answer in \textbf{standard form}, and be ready to
justify a sign to your partner.
\end{headlinebox}

\vspace{0.06in}

\begin{tcolorbox}[colback=white, colframe=black!40, title=\textbf{Tier R --- Remediate},
    fonttitle=\bfseries, arc=1.5mm, left=3mm, right=3mm, top=2mm, bottom=2mm, breakable]
\begin{enumerate}[leftmargin=1.7em, itemsep=9pt]
  \item \textbf{Add.} $(2x^2 + 5x - 3) + (4x^2 - x + 8) = $ \blank{5.5cm}

  \item \textbf{Subtract.} $(7x^2 - 4x + 1) - (3x^2 + 6x - 5)$

    \emph{Step 1 --- rewrite the second polynomial with every sign changed:}

    \vspace{2pt}
    $7x^2 - 4x + 1$ \; \blank{5.5cm}

    \vspace{3pt}
    \emph{Step 2 --- combine like terms:} \blank{5.5cm}

  \item \textbf{Distribute.} $4x(x^2 - 3x + 2) = $ \blank{5.5cm}

  \item \textbf{Multiply two binomials.} $(x + 6)(x - 2)$. Fill in the box, then collect.
    \begin{center}
    \renewcommand{\arraystretch}{1.9}
    \begin{tabular}{|c|c|c|}
    \hline
    $\times$ & $x$ & $-2$ \\ \hline
    $x$  & \blank{1.7cm} & \blank{1.7cm} \\ \hline
    $+6$ & \blank{1.7cm} & \blank{1.7cm} \\ \hline
    \end{tabular}
    \end{center}
    Answer in standard form: \blank{5.0cm}
\end{enumerate}
\end{tcolorbox}

\begin{tcolorbox}[colback=white, colframe=black!40, title=\textbf{Tier A --- Approaching Proficiency},
    fonttitle=\bfseries, arc=1.5mm, left=3mm, right=3mm, top=2mm, bottom=2mm, breakable]
\begin{enumerate}[leftmargin=1.7em, itemsep=10pt]
  \item \textbf{A vanishing degree.} $(5x^3 - 2x + 9) - (5x^3 + x^2 - 4x + 9) = $ \blank{5.0cm}

    \vspace{2pt}
    Both polynomials have degree $3$, but the answer does not. What is its degree? \blank{1.2cm}
    \; Explain what happened. \writeline

  \item \textbf{Binomial $\times$ trinomial.} $(3x - 2)(2x^2 - 5x + 4)$
    \begin{center}
    \renewcommand{\arraystretch}{1.9}
    \begin{tabular}{|c|c|c|c|}
    \hline
    $\times$ & $2x^2$ & $-5x$ & $+4$ \\ \hline
    $3x$ & \blank{1.7cm} & \blank{1.7cm} & \blank{1.7cm} \\ \hline
    $-2$ & \blank{1.7cm} & \blank{1.7cm} & \blank{1.7cm} \\ \hline
    \end{tabular}
    \end{center}
    Answer in standard form: \blank{6.0cm}

  \item \textbf{Special products, two variables.}

    \vspace{2pt}
    (a) $(4x + 3y)(4x - 3y) = $ \blank{4.5cm} \qquad (b) $(2x - 5y)^2 = $ \blank{5.0cm}

  \item \textbf{Predict, don't expand.} For the product $(2x^3 - 1)(-3x^4 + x)$:

    \vspace{2pt}
    Degree of the product: \blank{1.2cm} \quad Leading coefficient: \blank{1.4cm}

    \vspace{2pt}
    So its left arm points \blank{1.4cm} and its right arm points \blank{1.4cm}. How did you know
    without multiplying everything out? \writeline

  \item \textbf{Error analysis.} A classmate writes $(x + 4)^2 = x^2 + 16$.

    \vspace{2pt}
    (a) Test both sides at $x = 1$: \; left side $=$ \blank{1.4cm}, \; right side $=$ \blank{1.4cm}.

    \vspace{2pt}
    (b) Write the correct expansion: \blank{4.5cm}

    \vspace{2pt}
    (c) In one sentence, tell your classmate exactly what they left out and why. \writeline
\end{enumerate}
\end{tcolorbox}

\begin{tcolorbox}[colback=white, colframe=black!40, title=\textbf{Tier E --- Extension},
    fonttitle=\bfseries, arc=1.5mm, left=3mm, right=3mm, top=2mm, bottom=2mm, breakable]
\textbf{Part 1 --- The open box, expanded.} In Lesson 4.0 you read a table for the open-box volume
$V(x) = x(10-2x)(8-2x)$ but never saw its standard form. Build it now.
\begin{enumerate}[label=(\alph*), leftmargin=1.8em, itemsep=6pt]
  \item Multiply the two binomials first: $(10-2x)(8-2x) = $ \blank{5.0cm}
  \item Now multiply by $x$: \; $V(x) = $ \blank{5.5cm}
  \item Check against Lesson 4.0's table, which gave $V(1) = 48$. Substitute $x=1$ into your standard
        form: \blank{4.5cm} \; Does it match? \blank{1.4cm}
  \item Degree \blank{1.0cm}, leading coefficient \blank{1.0cm}. The polynomial itself is defined for
        \emph{every} real $x$, but the \emph{box} only exists for $0 < x < 4$. Why? \writeline
\end{enumerate}

\vspace{5pt}
\textbf{Part 2 --- Profit is a subtraction.} A small print shop sells $x$ hundred posters. Its revenue
and cost, in dollars, are
\[
  R(x) = 20x^2 + 150x \qquad\text{and}\qquad C(x) = 12x^2 + 90x + 400 .
\]
\begin{enumerate}[label=(\alph*), leftmargin=1.8em, itemsep=6pt]
  \item Profit is $P(x) = R(x) - C(x)$. Find $P(x)$ in standard form --- watch every sign.
        \blank{5.5cm}
  \item Evaluate $P(10)$: \blank{3.0cm} \; What does that number mean in context? \writeline
  \item The $-400$ in $P(x)$ came from the cost polynomial. What does it represent for the shop?
        \writeline
\end{enumerate}

\vspace{5pt}
\textbf{Part 3 --- Justify.} Two claims about degree. Decide whether each is \emph{always} true, and
support your decision with an example.
\begin{enumerate}[label=(\alph*), leftmargin=1.8em, itemsep=6pt]
  \item The \emph{product} of a degree-3 polynomial and a degree-4 polynomial always has degree 7.
        \writelines{2}
  \item The \emph{sum} of two degree-4 polynomials always has degree 4. \writelines{2}
\end{enumerate}
\end{tcolorbox}

\end{document}
